% !TEX root = ../Main.tex
\chapter{圆幂定理}

本章把和圆有关的几条乘积关系系统地串起来：弦切角、相交弦、割线、切割线，最后到托勒密。它们都靠\textbf{相似三角形}推出。

\section{弦切角定理}

过圆上一点 $P$ 作切线，再作弦 $PA$。则\textbf{弦切角}（切线与弦的夹角）等于它所夹的弧上的\textbf{圆周角}。

\begin{center}
\begin{tikzpicture}[>=Stealth,scale=1.0]
  \draw[thick] (0,0) circle (1.6);
  \coordinate (O) at (0,0);
  \coordinate (P) at (90:1.6);
  \coordinate (A) at (200:1.6);
  \coordinate (B) at (340:1.6);
  % tangent at P (horizontal)
  \draw[thick] (-1.7,1.6) -- (1.9,1.6);
  \draw (P) -- (A);
  \draw (P) -- (B);
  \draw[dashed] (A) -- (B);
  \fill (P) circle (1.2pt) node[above right]{$P$};
  \fill (A) circle (1.2pt) node[left]{$A$};
  \fill (B) circle (1.2pt) node[right]{$B$};
  \fill (O) circle (1pt) node[below]{$O$};
\end{tikzpicture}
\end{center}

即图中切线与弦 $PA$ 的夹角，等于弦 $PA$ 所对的圆周角 $\angle PBA$。这条是后面切割线定理的"零件"。

\section{相交弦定理}

\thm{相交弦定理}{
    两条弦 $AB$ 与 $CD$ 在圆内交于 $T$，则
    \[
    TA\cdot TB=TC\cdot TD.
    \]
}

\begin{center}
\begin{tikzpicture}[>=Stealth,scale=1.0]
  \draw[thick] (0,0) circle (1.6);
  \coordinate (C) at (150:1.6);
  \coordinate (A) at (30:1.6);
  \coordinate (B) at (210:1.6);
  \coordinate (D) at (320:1.6);
  \draw[name path=AB] (A) -- (B);
  \draw[name path=CD] (C) -- (D);
  \path[name intersections={of=AB and CD,by=T}];
  \fill (T) circle (1.2pt) node[above right]{$T$};
  \fill (C) circle (1.2pt) node[above left]{$C$};
  \fill (A) circle (1.2pt) node[above right]{$A$};
  \fill (B) circle (1.2pt) node[below left]{$B$};
  \fill (D) circle (1.2pt) node[below right]{$D$};
  \fill (0,0) circle (1pt) node[below left]{$O$};
\end{tikzpicture}
\end{center}

\pf[相似]{
    连 $AC,BD$（或 $AD,BC$）。由同弧所对圆周角相等及对顶角，$\triangle TAC\sim\triangle TDB$，得 $\dfrac{TA}{TD}=\dfrac{TC}{TB}$，即 $TA\cdot TB=TC\cdot TD$。
}

\section{割线定理}

\thm{割线定理}{
    $P$ 在圆外，过 $P$ 作两条割线，分别交圆于 $A,B$ 与 $C,D$，则
    \[
    AP\cdot PB=CP\cdot PD.
    \]
}

\begin{center}
\begin{tikzpicture}[>=Stealth,scale=1.0]
  \draw[thick] (0,0) circle (1.4);
  \coordinate (P) at (3,-0.3);
  \coordinate (B) at (140:1.4);
  \coordinate (A) at (38:1.4);
  \coordinate (D) at (220:1.4);
  \coordinate (C) at (-20:1.4);
  \draw (P) -- (B);
  \draw (P) -- (D);
  \fill (P) circle (1.2pt) node[right]{$P$};
  \fill (B) circle (1.2pt) node[above left]{$B$};
  \fill (A) circle (1.2pt) node[above]{$A$};
  \fill (D) circle (1.2pt) node[below left]{$D$};
  \fill (C) circle (1.2pt) node[below]{$C$};
  \fill (0,0) circle (1pt) node[left]{$O$};
\end{tikzpicture}
\end{center}

\section{切割线定理}

\thm{切割线定理}{
    $P$ 在圆外，$PA$ 切圆于 $A$，$PBC$ 是过 $P$ 的割线，则
    \[
    PA^2=PB\cdot PC.
    \]
}

\begin{center}
\begin{tikzpicture}[>=Stealth,scale=1.0]
  \draw[thick] (0,0) circle (1.4);
  \coordinate (P) at (2.9,-0.6);
  \coordinate (A) at (55:1.4);
  \coordinate (B) at (250:1.4);
  \coordinate (C) at (170:1.4);
  \draw (P) -- (A);
  \draw (P) -- (C);
  \fill (P) circle (1.2pt) node[right]{$P$};
  \fill (A) circle (1.2pt) node[above]{$A$};
  \fill (B) circle (1.2pt) node[below right]{$B$};
  \fill (C) circle (1.2pt) node[left]{$C$};
  \fill (0,0) circle (1pt) node[above left]{$O$};
\end{tikzpicture}
\end{center}

\pf[弦切角 + 相似]{
    连 $AB,AC$。由弦切角定理，$\angle PAB$ 等于弦 $AB$ 所对的圆周角 $\angle ACB=\angle ACP$；又 $\angle APB=\angle CPA$（公共角），故 $\triangle PAB\sim\triangle PCA$，得 $\dfrac{PA}{PC}=\dfrac{PB}{PA}$，即 $PA^2=PB\cdot PC$。
}

\section{托勒密定理}

\thm{托勒密定理}{
    圆内接四边形 $ABCD$ 中，
    \[
    AB\cdot CD+AD\cdot BC=AC\cdot BD.
    \]
    （两组对边乘积之和 $=$ 两条对角线乘积。）
}

\begin{center}
\begin{tikzpicture}[>=Stealth,scale=1.1]
  \draw[thick] (0,0) circle (1.5);
  \coordinate (A) at (110:1.5);
  \coordinate (B) at (200:1.5);
  \coordinate (C) at (320:1.5);
  \coordinate (D) at (40:1.5);
  \draw (A)--(B)--(C)--(D)--cycle;
  \draw (A)--(C);
  \draw (B)--(D);
  \fill (A) circle (1.2pt) node[above]{$A$};
  \fill (B) circle (1.2pt) node[left]{$B$};
  \fill (C) circle (1.2pt) node[below right]{$C$};
  \fill (D) circle (1.2pt) node[right]{$D$};
\end{tikzpicture}
\end{center}

\pf[构造相似]{
    在对角线 $AC$ 上取点 $E$，使 $\angle ABE=\angle DBC$。

    \textbf{一对相似}：$\triangle ABE\sim\triangle DBC$（$\angle ABE=\angle DBC$，且 $\angle BAE=\angle BAC=\angle BDC$ 为同弧圆周角），得
    \[
    \frac{AE}{DC}=\frac{AB}{DB}\ \Longrightarrow\ AB\cdot DC=AE\cdot DB.
    \]

    \textbf{另一对相似}：因 $\angle ABE=\angle DBC$，两边同加 $\angle EBD$ 得 $\angle ABD=\angle EBC$，又 $\angle ADB=\angle ACB=\angle ECB$（同弧），故 $\triangle ABD\sim\triangle EBC$，得
    \[
    \frac{AD}{EC}=\frac{BD}{BC}\ \Longrightarrow\ AD\cdot BC=EC\cdot BD.
    \]

    两式相加，并注意 $AE+EC=AC$：
    \[
    AB\cdot CD+AD\cdot BC=(AE+EC)\cdot BD=AC\cdot BD.
    \]
}
